\documentclass[conference]{IEEEtran}
\IEEEoverridecommandlockouts

\usepackage{amsmath,amssymb,amsfonts}
\usepackage{algorithm}
\usepackage{algpseudocode}
\usepackage{graphicx}
\usepackage{booktabs}
\usepackage{url}
\usepackage[caption=false,font=footnotesize]{subfig}

\DeclareMathOperator{\rank}{rank}
\DeclareMathOperator{\Tr}{Tr}
\DeclareMathOperator{\range}{range}
\newcommand{\bG}{\mathbf{G}}
\newcommand{\bS}{\mathbf{S}}
\newcommand{\bB}{\mathbf{B}}
\newcommand{\bW}{\mathbf{W}}
\newcommand{\bZ}{\mathbf{Z}}
\newcommand{\bH}{\mathbf{H}}
\newcommand{\bU}{\mathbf{U}}
\newcommand{\bJ}{\mathbf{J}}
\newcommand{\bD}{\mathbf{D}}
\newcommand{\bg}{\mathbf{g}}
\newcommand{\bv}{\mathbf{v}}
\newcommand{\bw}{\mathbf{w}}
\newcommand{\CN}{\mathcal{CN}}
\newcommand{\Hank}{\mathcal{H}}

\begin{document}

\title{Hankel-Structured Gerchberg--Saxton Channel\\Estimation for Rydberg Atomic MIMO Receivers}

\author{\IEEEauthorblockN{Author Name}
\IEEEauthorblockA{Department of Electrical and Electronics Engineering\\
Birla Institute of Technology and Science, Pilani\\
Email: f20220519@pilani.bits-pilani.ac.in}}

\maketitle

\renewcommand{\topfraction}{0.92}
\renewcommand{\bottomfraction}{0.75}
\renewcommand{\textfraction}{0.07}
\renewcommand{\floatpagefraction}{0.85}

\begin{abstract}
Rydberg atomic quantum receivers measure only the \emph{magnitude} of the
incident field superposed on a local reference, making multi-user channel
estimation a biased phase-retrieval rather than a linear-Gaussian problem. The
known solvers---biased Gerchberg--Saxton (GS) and its expectation--
maximisation refinement EM-GS---treat the channel as unstructured, and their
normalised mean-square error (NMSE) consequently does not improve as the
atomic array grows. We observe that on a uniform linear array each propagation
path contributes exactly one spatial complex exponential, so the Hankel
lifting of every channel column has rank at most the path count. Interleaving
a Cadzow low-rank Hankel projection with every exact EM-GS update yields HS-
GS; the model order comes from a held-out pilot residual, no oracle path count
is used, and the magnitude observation is never linearised. We also derive
unconstrained and geometry-constrained Cram\'er--Rao bounds for the exact
Rician likelihood. Over $2.8\times10^{4}$ trials, and \emph{under the tested
sparse geometric multipath model with i.i.d.\ uniform angles of arrival and
equal per-path power}, HS-GS improves on EM-GS by a mean of $2.85$~dB at
$N=32$, while EM-GS is flat in $N$ to $0.012$~dB. A controlled path-count
experiment confirms the mechanism: the gain decays monotonically from
$7.13$~dB at $L=2$ to zero at the Hankel rank cap. Robustness under clustered
propagation is not established and is the decisive next test.
\end{abstract}

\begin{IEEEkeywords}
Rydberg atomic receiver, channel estimation, phase retrieval, Hankel matrix,
Cadzow projection, constrained Cram\'er--Rao bound.
\end{IEEEkeywords}

%==========================================================================
\section{Introduction}
%==========================================================================
Rydberg atomic quantum receivers replace the RF front end with an ensemble of
highly excited atoms read out optically, giving sensitivity and bandwidth
flexibility unavailable to conventional antennas
\cite{cui2025,fancher2021,artusio2022}. Cui \emph{et al.} \cite{cui2025}
showed that the resulting atomic MIMO receiver does not obey the linear
transmission model: heterodyne readout against a local reference yields the
\emph{modulus} of the summed field, so recovery is a \emph{biased phase-
retrieval} (PR) problem \cite{netrapalli2015,dong2023}. They adapted the
Gerchberg--Saxton (GS) algorithm \cite{gerchberg1972} to this setting and
proposed EM-GS, which replaces the hard modulus substitution with the exact
posterior mean of the phase factor under the Rician law.

Both solvers treat $\bG\in\mathbb{C}^{N\times K}$ as $2NK$ free real
parameters. Under the magnitude model the Fisher information is block diagonal
across receive elements (Section~\ref{sec:crlb}), so estimation separates into
$N$ independent per-row problems: enlarging the array adds proportionally as
many unknowns as measurements and leaves the \emph{normalised} error
unchanged. Our simulations confirm this---EM-GS varies by only $0.012$~dB
across $N\in\{8,16,32\}$. Any improvement scaling with array size must
therefore come from structure coupling the rows.

Such structure exists whenever the array is a ULA and propagation is via a
small number of resolvable paths: each path imprints one spatial complex
exponential across the aperture. This is the sparsity that grid-based and
subspace estimators exploit in conventional MIMO, but it has not been used
inside a magnitude-only atomic receiver, where the non-linearity forbids the
usual linear least-squares machinery.

\textbf{Contributions.} We show that the Hankel lifting of each ULA channel
column has rank at most $L_k$, turning path sparsity into a grid-free low-rank
constraint requiring no angle estimation (Section~\ref{sec:hankel}); we
propose HS-GS, interleaving a Cadzow projection with every exact EM-GS update
at an order chosen from a held-out pilot residual rather than from oracle
knowledge (Section~\ref{sec:method}); we derive the unconstrained and
geometry-constrained CRLBs for this Rician likelihood, correcting a rank-2
misconstruction of the per-measurement Fisher information
(Section~\ref{sec:crlb}); and we isolate \emph{when} the method works with a
controlled path-count experiment, stating explicitly the model conditions
under which the claims hold (Sections~\ref{sec:results}--\ref{sec:limits}).

\textbf{Notation.} $(\cdot)^{T}$ and $(\cdot)^{H}$ are transpose and conjugate
transpose; $|\cdot|$ on a matrix is elementwise modulus; $\odot$ is the Hadamard
product; $\|\cdot\|_F$ is the Frobenius norm; $I_\nu$ is the order-$\nu$
modified Bessel function of the first kind.

%==========================================================================
\section{System Model}\label{sec:model}
%==========================================================================
\subsection{Magnitude-Only Measurement}
$K$ single-antenna users transmit to an $N$-element atomic ULA over $P$ pilot
symbols. Writing $\bG\in\mathbb{C}^{N\times K}$ for the channel,
$\bS\in\mathbb{C}^{K\times P}$ for the pilots, $\bB\in\mathbb{C}^{N\times P}$
for the local reference and $\bW\in\mathbb{C}^{N\times P}$ for noise, the
receiver observes
\begin{equation}\label{eq:obs}
\bZ=\bigl|\bG\bS+\bB+\bW\bigr|\in\mathbb{R}_{\ge0}^{N\times P},
\end{equation}
with $\bW$ i.i.d.\ $\CN(0,\sigma^2)$ added \emph{inside} the modulus. The phase
of $\bG\bS+\bB$ is unobservable; only the reference $\bB$ makes $\bG$
identifiable, through the cross term of
$|x+b|^2=|x|^2+2\Re\{xb^*\}+|b|^2$. Following \cite{cui2025}, the reference is
constant over $(n,p)$ with zero phase, the signal-to-noise ratio fixes
\begin{equation}\label{eq:snr}
\sigma^2=K/\mathrm{SNR}_{\mathrm{lin}},
\end{equation}
and the reference-to-signal ratio (RSR) is defined against a \emph{single}
user, giving $|b|=\sqrt{\mathrm{RSR}_{\mathrm{lin}}}$. Because $K$ users
transmit simultaneously, a nominal RSR of $12$~dB corresponds to $7.15$~dB
measured against the $K$-user total; all RSR values quoted here are the
single-user ones.

\subsection{Geometric ULA Channel}
The channel is generated as sparse geometric multipath. With half-wavelength
spacing the electrical angle of path $\ell$ of user $k$ is
$\psi_{\ell,k}=\pi\sin\theta_{\ell,k}$, and
\begin{equation}\label{eq:chan}
g_{n,k}=c\sum_{\ell=1}^{L_k}\alpha_{\ell,k}\,
e^{-\jmath(n-1)\psi_{\ell,k}},\qquad n=1,\dots,N,
\end{equation}
with atomic gain $c=1$ (so the electrical and atomic channels coincide),
$\alpha_{\ell,k}\sim\CN(0,\beta_k/L_k)$ and $\beta_k=1$.

The assumptions used in every experiment reported here are: sparse geometric
multipath with $L_k\sim\mathcal{U}\{3,\dots,7\}$ drawn i.i.d.\ per user per
realisation; angles of arrival i.i.d.\ \emph{uniform in $\theta$} (a
Kolmogorov--Smirnov test over $4465$ samples gives $p=0.150$ against uniform-
in-$\theta$, $p<10^{-4}$ against uniform-in-$\sin\theta$); equal average per-
path power, so no near--far effect; pilots i.i.d.\ $\CN(0,1)$, redrawn each
trial and rejection-sampled to full row rank with $P\ge2K$, hence \emph{not}
orthogonal (largest off-diagonal $|\bS\bS^{H}|/P$ is $0.233$); and flat,
narrowband propagation. The channel is \emph{not} a clustered 3GPP TR 38.901
or Saleh--Valenzuela model; Section~\ref{sec:limits} addresses the
consequences. Table~\ref{tab:params} collects the parameters.

\begin{table}[tb]
\caption{Simulation parameters. All experiments use the exact observation
\eqref{eq:obs}; the modulus is never linearised.}
\label{tab:params}
\centering
\footnotesize
\setlength{\tabcolsep}{4pt}
\begin{tabular}{@{}llcll@{}}
\toprule
\multicolumn{2}{@{}l}{\emph{Channel and array}} & &
\multicolumn{2}{l}{\emph{Estimators}} \\
\cmidrule(r){1-2}\cmidrule(l){4-5}
$N$ & $8,16,32$ & & $t_0$ & $50$ (GS, EM-GS) \\
$K$ & $3$ & & $T$ & $50$ (HS-GS) \\
$L_k$ & $\mathcal{U}\{3,\dots,7\}$ i.i.d. & & project every & $1$ iteration \\
$\theta_{\ell,k}$ & $\mathcal{U}[-\pi/2,\pi/2]$ & & $n_{\mathrm{cz}}$ & $4$ sweeps \\
$\alpha_{\ell,k}$ & $\CN(0,\beta_k/L_k)$, $\beta_k{=}1$ & & $\nu$ & $0.3$ held out \\
$c$ & $1$ & & $T_{\mathrm{sel}}$ & $20$ \\
spacing & $\lambda/2$ & & $\hat L$ grid & $1,\dots,r_{\max}$ \\
\cmidrule(r){1-2}\cmidrule(l){4-5}
\multicolumn{2}{@{}l}{\emph{Pilots and operating point}} & &
\multicolumn{2}{l}{\emph{Monte Carlo}} \\
\cmidrule(r){1-2}\cmidrule(l){4-5}
$P$ & $10,30$ (swept $6$--$40$) & & trials/pt & $400$ ($1200$ at $5$ pts) \\
$\bS$ & $\CN(0,1)$, redrawn & & total & $2.8\times10^{4}$ \\
& full row rank, $P\ge2K$ & & bootstrap & $2000$, paired \\
$\bB$ & constant, $\angle b=0$ & & & \\
SNR & $-5$ to $20$~dB, \eqref{eq:snr} & & & \\
RSR & $12$~dB (swept $0$--$24$) & & & \\
\bottomrule
\end{tabular}
\end{table}

\subsection{Baseline Solvers}
Both baselines exploit the fact that row $n$ of \eqref{eq:obs} depends only on
row $n$ of $\bG$, reducing \eqref{eq:obs} to $N$ independent biased PR problems
of dimension $K$. Biased GS \cite{cui2025} alternates a least-squares update
with the hard modulus substitution
$\hat{\mathbf{y}}=\bZ\odot e^{\jmath\angle(\bG\bS+\bB)}$. EM-GS instead uses
the exact conditional mean of the phase factor under the Rician law, which
introduces the Bessel ratio
\begin{equation}\label{eq:bessel}
R(\kappa)=\frac{I_1(\kappa)}{I_0(\kappa)},\qquad
\kappa=\frac{2}{\sigma^2}\,\bZ\odot\bigl|\bG\bS+\bB\bigr| ,
\end{equation}
so the restored observation is shrunk by $R(\kappa)\in[0,1)$ according to
measurement reliability. As $\kappa\to\infty$, $R\to1$ and EM-GS provably
reduces to biased GS---a limit reached in practice at high RSR
(Section~\ref{sec:results}-A).

%==========================================================================
\section{Hankel Structure of the ULA Channel}\label{sec:hankel}
%==========================================================================
Drop the user index and write one channel column as
$g_n=\sum_{\ell=1}^{L}\alpha_\ell z_\ell^{\,n}$ with
$z_\ell=e^{-\jmath\psi_\ell}$. Define the Hankel lifting with pencil parameter
$p$,
\begin{equation}\label{eq:lift}
\Hank(\bg)=
\begin{bmatrix}
g_0 & g_1 & \cdots & g_{p}\\
g_1 & g_2 & \cdots & g_{p+1}\\
\vdots & & \ddots & \vdots\\
g_{N-p-1} & \cdots & & g_{N-1}
\end{bmatrix}
\in\mathbb{C}^{(N-p)\times(p+1)} .
\end{equation}
Its $(q,r)$ entry is $g_{q+r}$, so
\begin{equation}\label{eq:sep}
\bigl[\Hank(\bg)\bigr]_{q,r}=\sum_{\ell=1}^{L}\alpha_\ell z_\ell^{\,q+r}
=\sum_{\ell=1}^{L}\alpha_\ell\,z_\ell^{\,q}\,z_\ell^{\,r},
\end{equation}
which separates in $q$ and $r$. With
$\bv_\ell=[1,z_\ell,\dots,z_\ell^{N-p-1}]^{T}$ and
$\bw_\ell=[1,z_\ell,\dots,z_\ell^{p}]^{T}$ this is exactly
\begin{equation}\label{eq:rank1}
\Hank(\bg)=\sum_{\ell=1}^{L}\alpha_\ell\,\bv_\ell\bw_\ell^{T},
\end{equation}
a sum of $L$ rank-one terms, hence
\begin{equation}\label{eq:rank}
\boxed{\ \rank\bigl(\Hank(\bg)\bigr)\le L.\ }
\end{equation}
\emph{Each propagation path creates one spatial complex exponential across the
ULA and therefore contributes exactly one rank-one component to the Hankel
lifting}, converting path sparsity into low-rank matrix structure without
discretising the angular domain. When the $z_\ell$ are distinct and $L$ does
not exceed the matrix dimensions, \eqref{eq:rank} holds with equality, and
conversely a sequence whose Hankel matrix has rank $L$ is a sum of $L$
exponentials \cite{cadzow1988}.

The constraint has a hard ceiling: since $\Hank(\bg)$ is $(N-p)\times(p+1)$
its rank cannot exceed $\min(N-p,p+1)$, maximised over $p$ at
\begin{equation}\label{eq:cap}
r_{\max}=\max_{1\le p\le N-1}\min(N-p,\,p+1)=\lceil N/2\rceil ,
\end{equation}
so $r_{\max}=4,8,16$ for $N=8,16,32$. If $L\ge r_{\max}$ then $\rank\le L$ is
satisfied by every vector and carries no information. This ceiling states
precisely \emph{when} the prior is informative, and
Section~\ref{sec:results}-C tests it directly.

%==========================================================================
\section{Proposed HS-GS Estimator}\label{sec:method}
%==========================================================================
\subsection{Cadzow Projection}
Let $\mathcal{S}_{\hat L}$ be the length-$N$ sequences whose Hankel lifting
has rank at most $\hat L$. Membership is the intersection of two sets with
exact projectors: fixed-rank matrices (truncated SVD, Frobenius-optimal by
Eckart--Young) and Hankel matrices (anti-diagonal averaging). Cadzow's algorithm \cite{cadzow1988} alternates them: build
$\bH\leftarrow\Hank(\bg)$; take its SVD
$\bH=\bU\boldsymbol{\Sigma}\mathbf{V}^{H}$; retain the $\hat L$ dominant
components as $\bH_{\mathrm{lr}}=\bU_{:,1:\hat L}\boldsymbol{\Sigma}_{1:\hat
L}\mathbf{V}_{:,1:\hat L}^{H}$; restore Hankel structure by averaging each
anti-diagonal, $\bg\leftarrow\Hank^{\dagger}(\bH_{\mathrm{lr}})$; and repeat
for $n_{\mathrm{cz}}=4$ sweeps.
The intersection is a non-convex variety and four sweeps do not reach its exact
nearest point, so $\mathcal{P}_{\hat L}$ is an \emph{approximate} projection
(Section~\ref{sec:results}-D quantifies the residual) and the iteration below
carries no convergence guarantee. When $\hat L\ge r_{\max}$ the routine returns
its input unchanged, so HS-GS reproduces EM-GS bit for bit.

\subsection{Interleaved Iteration and Order Selection}
The projection is applied \emph{inside} the exact iteration, not as post-
processing:
\begin{equation}\label{eq:map}
\bG^{(t)}=\mathcal{P}_{\hat L}\circ\mathcal{T}_{\mathrm{EM\text{-}GS}}
\bigl(\bG^{(t-1)}\bigr),\qquad t=1,\dots,T,
\end{equation}
with $\mathcal{P}_{\hat L}$ acting column-wise after \emph{every} one of the
$T=50$ updates (Algorithm~\ref{alg:hsgs}). This matters because EM-GS update
quality depends on the current iterate through $\kappa$ in \eqref{eq:bessel}:
a structurally cleaner iterate yields a better phase estimate, hence a better
next iterate.

HS-GS requires a target rank $\hat L$, and the true $L_k$ is unknown to any
practical receiver. The in-sample residual cannot supply it: constraining
$\bG$ to a smaller set can only worsen the fit to the pilots used for fitting,
so that criterion always prefers the largest candidate. Splitting the pilot
columns into a fitting set and a validation set of fraction $\nu=0.3$
penalises over-modelling correctly, and every quantity involved is observable:
\begin{equation}\label{eq:order}
\hat L=\arg\min_{1\le L\le r_{\max}}
\bigl\|\,\bZ_{\mathcal{V}}-\bigl|\bG_L\bS_{\mathcal{V}}+\bB_{\mathcal{V}}\bigr|\,\bigr\|_F^{2},
\end{equation}
where $\bG_L$ is the HS-GS estimate at order $L$ from the fitting columns
only, run for $T_{\mathrm{sel}}=20$ iterations. \textbf{No oracle rank is used
anywhere in this paper}; the selector returns a single scalar $\hat L$ applied
to all $K$ columns. The order matters in both directions: $\hat L<L$ truncates
genuine channel components, introducing a bias that does not vanish as noise
decreases, whereas $\hat L>L$ admits additional noise and weakens the prior.
Our selector systematically under-selects---with true orders $(7,4,6)$ a
typical choice is $\hat L=5$, and in Section~\ref{sec:results}-C the mean
$\hat L$ reaches only $11.55$ at $L=16$.

\begin{algorithm}[tb]
\caption{HS-GS}
\label{alg:hsgs}
\begin{algorithmic}[1]
\Require $\bZ,\bS,\bB,\sigma^2$; $T$, $n_{\mathrm{cz}}$, $\nu$, $T_{\mathrm{sel}}$
\State $\hat L\leftarrow$ held-out pilot selection \eqref{eq:order}
\State $\mathrm{active}\leftarrow(\hat L<r_{\max})$;\quad $\bG\leftarrow\varnothing$
\For{$t=1$ \textbf{to} $T$}
  \State $\bG\leftarrow\mathcal{T}_{\mathrm{EM\text{-}GS}}(\bG;\bZ,\bS,\bB,\sigma^2)$
         \Comment{exact update, eq.~\eqref{eq:bessel}}
  \If{$\mathrm{active}$}
    \For{$k=1$ \textbf{to} $K$}
      \State $\bG_{:,k}\leftarrow\mathcal{P}_{\hat L}(\bG_{:,k})$
             \Comment{$n_{\mathrm{cz}}$ Cadzow sweeps}
    \EndFor
  \EndIf
\EndFor
\State \Return $\hat\bG_{\mathrm{HS}}=\bG$
\end{algorithmic}
\end{algorithm}

\subsection{Complexity}
One EM-GS update costs $O(NPK)$ once the $K\times K$ Gram matrix is cached;
the projection adds $Kn_{\mathrm{cz}}$ SVDs of an $(N-p)\times(p+1)$ matrix
per iteration, i.e.\ $O(Kn_{\mathrm{cz}}N^{3})$, and order selection
multiplies the estimation by the $r_{\max}$ candidates. Against a
\emph{chained} EM-GS baseline---the correct reference, since HS-GS issues $T$
single-iteration calls costing $1.005$--$1.038\times$ in overhead alone---HS-
GS at fixed order costs $1.49\times$, $1.32\times$ and $1.22\times$ at
$N=8,16,32$, the ratio \emph{falling} with $N$ because the exact step grows
faster than the projection at these dimensions. Order selection dominates the
full estimator, at $57.7\%$, $74.1\%$ and $85.5\%$ of its runtime; shrinking
the candidate set or reusing iterates across candidates is therefore the
obvious practical target.

%==========================================================================
\section{Performance Benchmark}\label{sec:crlb}
%==========================================================================
\subsection{Likelihood and Fisher Information}
Conditional on $\bG,\bS,\bB$, the noiseless field at $(n,p)$ is a known complex
constant $\lambda=(\bG\bS+\bB)_{n,p}$ and the measurement $z=|\lambda+w|$ is
Rice distributed:
\begin{equation}\label{eq:rice}
p(z\mid\lambda)=\frac{2z}{\sigma^2}
\exp\!\Bigl(-\frac{z^2+|\lambda|^2}{\sigma^2}\Bigr)
I_0\!\Bigl(\frac{2z|\lambda|}{\sigma^2}\Bigr).
\end{equation}
Crucially \eqref{eq:rice} depends on $\lambda$ \emph{only through
$a=|\lambda|$}: the phase does not appear. Differentiating and using
$\frac{d}{dx}\log I_0(x)=R(x)$,
\begin{equation}\label{eq:score}
\frac{\partial\log p}{\partial a}=\frac{2}{\sigma^2}\bigl[zR(\kappa)-a\bigr],
\qquad \kappa=\frac{2za}{\sigma^2},
\end{equation}
so a single measurement carries the scalar information
$\beta(a,\sigma^2)=\mathbb{E}\bigl[(\partial_a\log p)^2\bigr]$, and its
contribution to the Fisher matrix in the real coordinates
$\boldsymbol{\eta}=[\Re\mathrm{vec}\,\bG;\Im\mathrm{vec}\,\bG]\in\mathbb{R}^{2NK}$
is
\begin{equation}\label{eq:fisher}
\bJ_{n}=\sum_{p=1}^{P}4\beta(a_{n,p},\sigma^2)\,
\mathbf{u}_{n,p}\mathbf{u}_{n,p}^{T},
\end{equation}
with $\mathbf{u}_{n,p}=\partial a_{n,p}/\partial\boldsymbol{\eta}_n$.
\textbf{Each magnitude measurement therefore contributes information of rank
one, not rank two.} Crediting both quadratures overstates the information and
yields a bound too low by $0.17$--$4.64$~dB across our $36$ operating points.
An identifiability check settles the form: with $K=3$ (six real unknowns) and
$P=3$ measurements the per-row problem is unidentifiable and \eqref{eq:fisher}
is correctly singular, whereas the rank-two construction claims full rank.
Fig.~\ref{fig:snr} corroborates it---EM-GS sits within $0.05$~dB of the rank-
one bound at high SNR, and would appear to violate the rank-two one. Because
row $n$ of $\bG\bS$ enters only measurement row $n$, $\bJ$ is block diagonal
across receive elements, which is why the unconstrained bound---and any
estimator attaining it---is flat in $N$ under a normalised metric.

\subsection{Constrained Bound on the Geometric Manifold}
The unconstrained bound governs estimators that treat $\bG$ as $2NK$ free
parameters. HS-GS instead exploits \eqref{eq:chan}, whose parameters are
$\boldsymbol{\phi}=\{\psi_{\ell,k},\Re\alpha_{\ell,k},\Im\alpha_{\ell,k}\}
\in\mathbb{R}^{3\sum_k L_k}$---about $45$ real numbers, \emph{independent of
$N$}. With $\bD=\partial\boldsymbol{\eta}/\partial\boldsymbol{\phi}$ the
textbook form $\bD(\bD^{T}\bJ\bD)^{-1}\bD^{T}$ is unusable here, because $\bD$
is rank deficient: at $N=8$ the mean numerical rank of $\bD$ is $40.4$ against
$3\sum_k L_k\approx44.7$ structural parameters, and $28.1\%$ of realisations
have $3\sum_k L_k>2NK=48$ outright, so $\boldsymbol{\phi}\mapsto\bG$ is not
injective. We therefore use the Gorman--Hero / Stoica--Ng tangent-space form
\cite{gorman1990,stoica1998}, with $\bU$ an orthonormal basis of $\range(\bD)$
obtained by SVD at relative tolerance $10^{-9}$:
\begin{equation}\label{eq:ccrb}
\mathrm{CCRB}=\bU\bigl(\bU^{T}\bJ\bU\bigr)^{-1}\bU^{T},
\end{equation}
which depends on the tangent \emph{space} and is thus invariant to
over-parameterisation. The worst conditioning of $\bU^{T}\bJ\bU$ over all
operating points is then $80.1$, at the degenerate point $P=6$; elsewhere it is
below $35$. Both bounds are averaged over the same $400$ channel realisations
used for the estimator curves, so the comparison is paired on the channel draw.

\subsection{Bias Caveat}
Both bounds constrain the covariance of \emph{unbiased} estimators, whereas
\begin{equation}\label{eq:mse}
\mathrm{MSE}=\mathrm{variance}+\mathrm{bias}^2 .
\end{equation}
All three estimators run to a fixed budget with the shrinking update
\eqref{eq:bessel}, and HS-GS additionally truncates rank, removing signal
energy; none is unbiased, so neither bound applies strictly. Two consequences
govern how the bounds are read. The unconstrained CRLB is the appropriate
reference for GS and EM-GS only: HS-GS uses prior information that bound does
not model, so it may legitimately fall below it and must be judged against the
CCRB. And the CCRB, itself an unbiased-estimation bound on the manifold, does
not forbid a sufficiently biased structured estimator from falling below it
either. Section~\ref{sec:results}-E reports both effects with the accompanying
shrinkage measurements.

%==========================================================================
\section{Simulation Setup}\label{sec:setup}
%==========================================================================
All results use the exact observation \eqref{eq:obs}; the modulus is never
linearised, which is verified per trial. NMSE is pooled as a ratio of sums,
\begin{equation}\label{eq:nmse}
\mathrm{NMSE}_G=10\log_{10}
\frac{\sum_{t}\|\hat\bG_t-\bG_t\|_F^2}{\sum_{t}\|\bG_t\|_F^2},
\end{equation}
never as a mean of per-trial decibels: \eqref{eq:nmse} estimates the ensemble
quantity a CRLB bounds, whereas a mean-of-decibels is strictly smaller by
Jensen's inequality. The factor is $10$, not $20$, since the arguments are
energies. Within an operating point all estimators see bit-identical
realisations, so the reported gain
\begin{equation}\label{eq:gain}
\Delta_{\mathrm{HS}}=10\log_{10}
\frac{\sum_t\|\hat\bG^{\mathrm{EM}}_t-\bG_t\|_F^2}
     {\sum_t\|\hat\bG^{\mathrm{HS}}_t-\bG_t\|_F^2}
\end{equation}
and its $95\%$ paired bootstrap interval ($2000$ resamples) benefit from
cancellation of the common channel draw; realisations are independent
\emph{across} operating points. Unless noted, $400$ trials per point are used,
extended to $1200$ at the five points whose intervals initially straddled
zero, for $2.8\times10^{4}$ trials in total.

%==========================================================================
\section{Results and Discussion}\label{sec:results}
%==========================================================================

\begin{figure}[tb]
\centering
\includegraphics[width=0.95\columnwidth]{fig/fig1_nmse_vs_snr.pdf}
\caption{NMSE versus SNR at $N=32$, $P=30$, RSR $=12$~dB. EM-GS is
indistinguishable from the unconstrained CRLB across the range, so the bound is
tight for the unstructured problem. HS-GS lies between the two bounds,
recovering part of the structural headroom identified by the geometric CCRB.}
\label{fig:snr}
\end{figure}

\begin{figure}[tb]
\centering
\includegraphics[width=0.92\columnwidth]{fig/fig2_gain_vs_N.pdf}
\caption{Paired HS-GS gain over EM-GS versus array size. Faint markers are the
six individual SNR operating points; solid lines are their mean. The gain grows
monotonically with $N$; at $N=8$ the individual points straddle zero, so no
single scalar summarises that array size.}
\label{fig:scaling}
\end{figure}

\subsection{Main Result and Array-Size Scaling} We first fix the operating
point. EM-GS improves on biased GS by $1.31$~dB at RSR $=0$~dB, but the
advantage collapses to $0.113$~dB at $12$~dB and $0.002$~dB at $24$~dB as
$R(\kappa)\to1$ and the two algorithms provably coincide. RSR $=12$~dB,
following the physical argument of \cite{cui2025}, keeps the baselines
separable and implies that the gains below are attributable to the structural
projection rather than to the choice of exact step. Fig.~\ref{fig:snr} shows
the principal comparison at $N=32$, $P=30$. EM-GS
tracks the unconstrained CRLB to within $0.05$~dB for SNR $\ge5$~dB, confirming
that it is close to efficient for the unstructured problem and that the
corrected bound is credible. The geometric CCRB lies $0.87$--$9.98$~dB below the
unconstrained bound across all operating points, quantifying how much the
structure is worth in principle. HS-GS occupies part of that gap, improving on
EM-GS by $2.27$~dB at $5$~dB SNR and by $2.27$--$3.04$~dB across the sweep.

Fig.~\ref{fig:scaling} isolates the scaling. Averaged over the twelve
operating points at each array size, the gain is $-0.19$, $+0.78$ and
$+2.85$~dB for $N=8,16,32$, the per-trial win rate $33.5\%$, $74.0\%$ and
$95.3\%$, and the constraint active in $57.5\%$, $92.7\%$ and $99.6\%$ of
trials; the largest single-point gain is $4.33$~dB ($N=32$, $P=10$, SNR
$=-5$~dB). Under the same averaging EM-GS varies by only $0.012$~dB across the
three array sizes ($-8.846$, $-8.848$, $-8.836$~dB), exactly as the block-
diagonal Fisher structure of Section~\ref{sec:crlb} requires. The
$N$-dependence therefore belongs to the structural prior, not to a degrading
baseline.

The mechanism is the ratio of path count to rank ceiling \eqref{eq:cap}. At
$N=8$, $r_{\max}=4$ while $L_k\in\{3,\dots,7\}$, so the selected order often
reaches the cap, the projection is inactive in $43\%$ of trials, and where
active it truncates real components. At $N=32$, $r_{\max}=16\gg L_k$ and the
constraint is active in $99.6\%$ of trials. We note that at $N=8$ the sign of
any scalar summary depends on the pooling---the per-point mean is $-0.19$~dB
but energy-pooling all trials gives $+0.81$~dB---because HS-GS wins at low SNR
and loses at high SNR there. That high-SNR loss is a bias-dominated floor: the
truncation bias does not shrink with noise, so as variance falls the trade
stops paying. At $N=16$ and $32$ every pooling agrees in sign.

\begin{figure}[tb]
\centering
\includegraphics[width=0.97\columnwidth]{fig/fig3_pathcount.pdf}
\caption{Controlled path count ($N=32$, $P=30$, SNR $=5$~dB, RSR $=12$~dB,
$400$ trials per point, $L$ fixed and identical across users). (a) The two
baselines are flat in $L$---path count does not affect an estimator that
ignores structure---which confirms the sweep is not confounded; HS-GS rises to
meet them. (b) The paired gain decays monotonically to zero as $L$ approaches the rank
ceiling $r_{\max}=\lceil N/2\rceil=16$; the $95\%$ paired bootstrap band is
plotted but is narrower than the markers. The hypothesis producing this figure
was recorded before the experiment was run.}
\label{fig:pathcount}
\end{figure}

\subsection{Controlled Path Count: When and Why the Method Works}
Every experiment above draws $L_k$ at random, so none isolates the path count.
Fig.~\ref{fig:pathcount} fixes $L_k=L$ identically across users and sweeps it to
$r_{\max}$. The gain decays strictly monotonically,
\[
7.13,\;3.58,\;1.84,\;1.03,\;0.59,\;0.27,\;0.04,\;-0.12~\text{dB}
\]
for $L=2,4,\dots,16$, with the per-trial win rate falling from $100\%$ to
$45.3\%$ and EM-GS flat to $0.140$~dB across the sweep. The interpretation is
direct: when $L\ll r_{\max}$ the Hankel lifting carries substantial redundant
structure, and projecting onto the low-rank Hankel set suppresses estimation
error that does not share that structure. As $L\to r_{\max}$ the prior admits
almost every vector and the benefit disappears.

Two details deserve emphasis. First, the null is reached at $L=14$ (gain
$+0.04$~dB, CI $[-0.05,+0.12]$), one grid point \emph{before} the ceiling, and
at $L=16$ the gain is small but significantly negative ($-0.12$~dB, CI
$[-0.195,-0.041]$). This is explained by the selector rather than by the theory:
mean $\hat L$ reaches only $11.55$ at $L=16$, so the constraint is still active
in $90.2\%$ of trials and what it does there is discard genuine channel
components---the $\hat L<L$ regime of Section~\ref{sec:method}-B. Second, this
result rules out the most plausible rival explanation. If the improvement were
generic denoising it would persist at large $L$, where the estimate is equally
noisy; instead it vanishes and reverses exactly where the structural prior
becomes vacuous.

\subsection{Structural Diagnostic}
Fig.~\ref{fig:spectrum} shows why the prior exists and what the projection
achieves. The Hankel spectrum of the true channel falls to $\sim\!10^{-16}$
immediately after index $L$: the low-rank property is exact, as
\eqref{eq:rank} states. The EM-GS estimate has no such cliff---noise fills the
whole spectrum, so the estimate is not a sum of $L$ exponentials and the
structural information is absent from it. After Cadzow a cliff reappears at
the correct index but descends only to $\sim\!10^{-3}$. That shortfall is the
four-sweep projection being approximate, and it upper-bounds how much
structure can be imposed---consistent with HS-GS attaining only part of the
CCRB headroom.

\begin{figure}[tb]
\centering
\includegraphics[width=0.95\columnwidth]{fig/fig4_hankel_spectrum.pdf}
\caption{Normalised Hankel singular values for one channel column, $N=32$,
$L=3$, SNR $=5$~dB. The true channel is exactly rank $L$; the EM-GS estimate is
full rank; Cadzow restores the cliff approximately, to $\sim\!10^{-3}$ rather
than to machine precision.}
\label{fig:spectrum}
\end{figure}

\subsection{Distance to the Bounds, and Crossings}
HS-GS lies above the geometric CCRB at $34$ of $36$ operating points, by
$0.32$--$8.85$~dB, and the gap \emph{grows} with $N$ ($+0.32$ to $+1.41$~dB at
$N=8$; $+4.29$ to $+4.82$~dB at $N=32$, $P=30$). HS-GS therefore captures more
structural gain in absolute terms as the array grows but a smaller share of what
is available---the remaining headroom lies in the two weaknesses identified
above, the shared scalar $\hat L$ and the approximate projection.

Curves do fall below the plotted bounds, and two situations must be separated.
HS-GS lies below the \emph{unconstrained} CRLB at $26$ of $36$ points, by up
to $6.94$~dB, but this is not a violation: that bound governs estimators
treating $\bG$ as $2NK$ free parameters, whereas HS-GS exploits a prior it
does not model. The CCRB is HS-GS's appropriate bound, and it is respected at
$34$ of $36$ points. The remaining crossings are consistent with bias, as
\eqref{eq:mse} permits. GS and EM-GS fall below the unconstrained CRLB at $18$
and $19$ of $36$ points, but the excursions are strongly SNR-ordered: up to
$2.34$ and $2.62$~dB at $-5$~dB, $1.05$ and $1.16$~dB at $0$~dB, and at most
$0.05$~dB for SNR $\ge15$~dB, within Monte Carlo error. The measured shrinkage
coefficient $c=\langle\hat\bG,\bG\rangle/\|\bG\|_F^{2}$ follows the same
ordering: at $(N,P)=(32,30)$, $c_{\mathrm{EM}}=0.944$ at $-5$~dB but $1.002$
at $+5$~dB, where the crossings have vanished. HS-GS falls below the CCRB at
exactly the two lowest SNRs of $(N,P)=(8,10)$, precisely the points of
strongest shrinkage: $c_{\mathrm{HS}}=0.561$ at $-5$~dB and $0.776$ at $0$~dB,
against $0.976$ at $20$~dB where no crossing occurs. HS-GS shrinks more than
the baselines at every SNR, since rank truncation discards signal energy on
top of the shrinkage already in \eqref{eq:bessel}. We do not claim bias
mathematically accounts for every crossing, only that the shrinkage
measurements track them in magnitude and location and that no other mechanism
was identified.

%==========================================================================
\section{Limitations and Future Work}\label{sec:limits}
%==========================================================================
All results here were produced under the i.i.d.\ uniform-AoA, equal-per-path-
power, flat geometric model of Section~\ref{sec:model}, so the defensible
claim is model-conditional: \emph{under the tested sparse geometric multipath
model, exploiting the low-rank Hankel structure of the ULA response improves
channel-estimation NMSE relative to EM-GS, by a mean of $2.85$~dB at $N=32$.}
No result here establishes robustness under clustered propagation. The
decisive next experiment is to repeat the principal comparison under a
clustered 3GPP TR 38.901 generator \cite{tr38901}---a falsification test, not
a confirmation exercise. Under clustering, rays within a cluster occupy a
narrow angular window, so their $z_\ell$ are tightly packed and $\Hank(\bg)$
becomes only \emph{numerically} low rank with a rapidly decaying spectrum
rather than an exact cliff. Hard truncation, exactly right for the latter, is
then a bias-inducing approximation, and the order selector may either adapt to
the effective rank or chase the decaying tail. Since the $N=8$ high-SNR
degradation already traces to truncation bias, clustering could plausibly
weaken or eliminate the gain.

Three further limitations are noted. A single scalar $\hat L$ is imposed on
all $K$ users despite their differing true orders; per-user selection would
change only the selector and would isolate how much of the CCRB gap is order
misspecification. The iteration \eqref{eq:map} has no convergence guarantee---
non-convex constraint set, approximate projection---and is defined by a fixed
budget. Finally, all users have equal average power, so near--far effects are
untested.

%==========================================================================
\section{Conclusion}\label{sec:conclusion}
%==========================================================================
The ULA response of a sparse geometric channel makes the Hankel lifting of
each channel column rank-deficient, and this structure can be enforced inside
an exact biased phase-retrieval iteration without linearising the magnitude
observation and without estimating any angle. The resulting HS-GS estimator
improves on EM-GS by a mean of $2.85$~dB at $N=32$ under the tested model,
whereas the unstructured baselines are flat in array size as the block-
diagonal Fisher information requires. A controlled path-count experiment both
confirms the mechanism and delimits it: the gain scales with the redundancy
$r_{\max}/L$ and vanishes when the path count reaches the Hankel ceiling.
Against a constrained CRLB derived for this likelihood, HS-GS attains part but
not all of the available structural gain, localising the remaining headroom in
per-user order selection and in the approximate projection. Validating the
method under clustered propagation remains the essential next step.

%==========================================================================
\begin{thebibliography}{99}
\footnotesize

\bibitem{cui2025}
M.~Cui, Q.~Zeng, and K.~Huang, ``Towards atomic MIMO receivers,''
\emph{IEEE J. Sel. Areas Commun.}, vol.~43, no.~3, pp.~659--673, Mar. 2025.

\bibitem{fancher2021}
C.~T. Fancher, D.~R. Scherer, M.~C. S.~John, and B.~L. S.~Marlow, ``Rydberg atom
electric field sensors for communications and sensing,'' \emph{IEEE Trans.
Quantum Eng.}, vol.~2, pp.~1--13, 2021.

\bibitem{artusio2022}
A.~Artusio-Glimpse, M.~T. Simons, N.~Prajapati, and C.~L. Holloway, ``Modern RF
measurements with hot atoms: A technology review of Rydberg atom-based radio
frequency field sensors,'' \emph{IEEE Microw. Mag.}, vol.~23, no.~5,
pp.~44--56, May 2022.

\bibitem{netrapalli2015}
P.~Netrapalli, P.~Jain, and S.~Sanghavi, ``Phase retrieval using alternating
minimization,'' \emph{IEEE Trans. Signal Process.}, vol.~63, no.~18,
pp.~4814--4826, Sep. 2015.

\bibitem{dong2023}
J.~Dong, L.~Valzania, A.~Maillard, T.-A. Pham, S.~Gigan, and M.~Unser, ``Phase
retrieval: From computational imaging to machine learning: A tutorial,''
\emph{IEEE Signal Process. Mag.}, vol.~40, no.~1, pp.~45--57, Jan. 2023.

\bibitem{gerchberg1972}
R.~W. Gerchberg and W.~O. Saxton, ``A practical algorithm for the determination
of phase from image and diffraction plane pictures,'' \emph{Optik}, vol.~35,
no.~2, pp.~237--246, 1972.

\bibitem{cadzow1988}
J.~A. Cadzow, ``Signal enhancement---A composite property mapping algorithm,''
\emph{IEEE Trans. Acoust., Speech, Signal Process.}, vol.~36, no.~1,
pp.~49--62, Jan. 1988.

\bibitem{gorman1990}
J.~D. Gorman and A.~O. Hero, ``Lower bounds for parametric estimation with
constraints,'' \emph{IEEE Trans. Inf. Theory}, vol.~36, no.~6, pp.~1285--1301,
Nov. 1990.

\bibitem{stoica1998}
P.~Stoica and B.~C. Ng, ``On the Cram\'er--Rao bound under parametric
constraints,'' \emph{IEEE Signal Process. Lett.}, vol.~5, no.~7, pp.~177--179,
Jul. 1998.

\bibitem{tr38901}
3GPP, ``Study on channel model for frequencies from 0.5 to 100 GHz,''
3rd Generation Partnership Project, Tech. Rep. TR 38.901, V16.1.0, 2020.

\end{thebibliography}

\end{document}
